\BeginAppendices

% ==========================================
% Приложение А - Код Gemini
% ==========================================
\prilojenie{Листинг скрипта, сгенерированного Gemini 3 Pro}
\label{app:gemini_code}

\inputcode[style=pythonstyle, caption={Скрипт проверки PDF (Gemini 3 Pro)}]{../code/Gemini3_check_pdf.py}

% ==========================================
% Приложение Б - Код Qwen
% ==========================================
\prilojenie{Листинг скрипта, сгенерированного Qwen 3 Coder}
\label{app:qwen_code}

\inputcode[style=pythonstyle, caption={Скрипт проверки PDF (Qwen 3 Coder)}]{../code/Qwen3_code_check_pdf.py}

% ==========================================
% Приложение В - Код YandexGPT
% ==========================================
\prilojenie{Листинг скрипта, сгенерированного YandexGPT 5.1}
\label{app:yagpt_code}

\inputcode[style=pythonstyle, caption={Скрипт проверки PDF (YandexGPT 5.1)}]{../code/YaGPT5_1_check_pdf.py}

% ==========================================
% Приложение Г - Лог Gemini
% ==========================================
\prilojenie{Протокол взаимодействия с Gemini 3 Pro}
\label{app:gemini_log}

% Разбивка на части для экономии памяти TeX
\inputcode[style=logstyle, lastline=291, caption={Лог Gemini 3 Pro (Часть 1)}]{../code/results/gemini_think_high/gemini_main_res.txt}
\clearpage
\inputcode[style=logstyle, firstline=292, caption={Лог Gemini 3 Pro (Часть 2)}]{../code/results/gemini_think_high/gemini_main_res.txt}

% ==========================================
% Приложение Д - Лог Qwen
% ==========================================
\prilojenie{Протокол взаимодействия с Qwen 3 Coder}
\label{app:qwen_log}

\inputcode[style=logstyle, lastline=289, caption={Лог Qwen 3 Coder (Часть 1)}]{../code/results/qwen/qwen_main_res.txt}
\clearpage
\inputcode[style=logstyle, firstline=290, caption={Лог Qwen 3 Coder (Часть 2)}]{../code/results/qwen/qwen_main_res.txt}

% ==========================================
% Приложение Е - Лог YandexGPT
% ==========================================
\prilojenie{Протокол взаимодействия с YandexGPT 5.1 Pro}
\label{app:yagpt_log}

\inputcode[style=logstyle, lastline=312, caption={Лог YandexGPT 5.1 Pro (Часть 1)}]{../code/results/ya_think/yagpt_main_res.txt}
\clearpage
\inputcode[style=logstyle, firstline=314, caption={Лог YandexGPT 5.1 Pro (Часть 2)}]{../code/results/ya_think/yagpt_main_res.txt}

% ==========================================
% Приложение Ж - Мульти-запуск Gemini High
% ==========================================
\prilojenie{Результаты контрольных запусков Gemini 3 Pro (Thinking High)}
\label{app:gemini_high_multirun}

\inputcode[style=logstyle, caption={Gemini 3 Pro (High) --- Запуск 1}]{../code/results/gemini_think_high/1_res.txt}
\clearpage
\inputcode[style=logstyle, caption={Gemini 3 Pro (High) --- Запуск 2}]{../code/results/gemini_think_high/2_res.txt}
\clearpage
\inputcode[style=logstyle, caption={Gemini 3 Pro (High) --- Запуск 3}]{../code/results/gemini_think_high/3_res.txt}
\clearpage
\inputcode[style=logstyle, caption={Gemini 3 Pro (High) --- Запуск 4}]{../code/results/gemini_think_high/4_res.txt}

% ==========================================
% Приложение И - Мульти-запуск Gemini Low
% ==========================================
\prilojenie{Результаты контрольных запусков Gemini 3 Pro (Thinking Low)}
\label{app:gemini_low_multirun}

\inputcode[style=logstyle, caption={Gemini 3 Pro (Low) --- Запуск 1}]{../code/results/gemini_think_low/1_res.txt}
\clearpage
\inputcode[style=logstyle, caption={Gemini 3 Pro (Low) --- Запуск 2}]{../code/results/gemini_think_low/2_res.txt}
\clearpage
\inputcode[style=logstyle, caption={Gemini 3 Pro (Low) --- Запуск 3}]{../code/results/gemini_think_low/3_res.txt}
\clearpage
\inputcode[style=logstyle, caption={Gemini 3 Pro (Low) --- Запуск 4}]{../code/results/gemini_think_low/4_res.txt}

% ==========================================
% Приложение К - Мульти-запуск YaGPT (Think)
% ==========================================
\prilojenie{Результаты контрольных запусков YandexGPT 5.1 (Thinking On)}
\label{app:yagpt_think_multirun}

\inputcode[style=logstyle, caption={YandexGPT 5.1 (Thinking On) --- Запуск 1}]{../code/results/ya_think/1_res.txt}
\clearpage
\inputcode[style=logstyle, caption={YandexGPT 5.1 (Thinking On) --- Запуск 2}]{../code/results/ya_think/2_res.txt}
\clearpage
\inputcode[style=logstyle, caption={YandexGPT 5.1 (Thinking On) --- Запуск 3}]{../code/results/ya_think/3_res.txt}
\clearpage
\inputcode[style=logstyle, caption={YandexGPT 5.1 (Thinking On) --- Запуск 4}]{../code/results/ya_think/4_res.txt}

% ==========================================
% Приложение Л - Мульти-запуск YaGPT (No Think)
% ==========================================
\prilojenie{Результаты контрольных запусков YandexGPT 5.1 (Thinking Off)}
\label{app:yagpt_no_think_multirun}

\inputcode[style=logstyle, caption={YandexGPT 5.1 (Thinking Off) --- Запуск 1}]{../code/results/ya_no_think/1_res.txt}
\clearpage
\inputcode[style=logstyle, caption={YandexGPT 5.1 (Thinking Off) --- Запуск 2}]{../code/results/ya_no_think/2_res.txt}
\clearpage
\inputcode[style=logstyle, caption={YandexGPT 5.1 (Thinking Off) --- Запуск 3}]{../code/results/ya_no_think/3_res.txt}
\clearpage
\inputcode[style=logstyle, caption={YandexGPT 5.1 (Thinking Off) --- Запуск 4}]{../code/results/ya_no_think/4_res.txt}

% ==========================================
% Приложение М - Мульти-запуск Qwen
% ==========================================
\prilojenie{Результаты контрольных запусков Qwen 3 Coder}
\label{app:qwen_multirun}

\inputcode[style=logstyle, caption={Qwen 3 Coder --- Запуск 1}]{../code/results/qwen/1_res.txt}
\clearpage
\inputcode[style=logstyle, caption={Qwen 3 Coder --- Запуск 2}]{../code/results/qwen/2_res.txt}
\clearpage
\inputcode[style=logstyle, caption={Qwen 3 Coder --- Запуск 3}]{../code/results/qwen/3_res.txt}
\clearpage
\inputcode[style=logstyle, caption={Qwen 3 Coder --- Запуск 4}]{../code/results/qwen/4_res.txt}